\subsection{Example 6: Platform Dashboard --- Internal Performance Overview}
\label{subsec:example-platform-dashboard}

The first example exercises the platform-wide analytics surface that an
internal operator consults to monitor aggregate health. Figure~\ref{fig:example-7}
shows the dashboard rendered for a seven-day window.

\begin{figure}[H]
    \centering
    \includegraphics[width=\textwidth]{images/illustrative_example_7.png}
    \caption{Platform Dashboard for a seven-day window.}
    \label{fig:example-7}
\end{figure}

When the page mounts, the React \texttt{DashboardPage} component issues a
single request to \texttt{GET /analytics/summary?days=7}. The backend
\texttt{analytics\_summary} handler executes three filtered counts against
the \texttt{events} table---one per event type within the time window---and
derives the click-through rate as the ratio of clicks to impressions.
Estimated revenue is computed as $\text{impressions} \times
\overline{\text{bid\_cpm}} / 1000$, where the average CPM is taken across
all currently active ads. A fourth query returns the ten ads with the
highest impression counts in the window, joined to the \texttt{ads} table
to recover product names. The response is a single \texttt{AnalyticsSummary}
object containing the five scalar metrics and the \texttt{top\_ads} array.

The frontend renders the five scalars as KPI cards and the array as a
sortable table. The CTR column carries an inline colour indicator that
encodes performance bands at a glance: green above three percent, amber
between one and three percent, red below one percent. The time-range
selector in the header re-issues the request with a different \texttt{days}
parameter; no client-side aggregation is required.

% ============================================================

\subsection{Example 7: Advertiser Rollup --- Multi-Campaign Budget Health}
\label{subsec:example-advertiser-rollup}

The second example shows the per-advertiser view that a self-serve client
consults to track spend across all of their campaigns. Figure~\ref{fig:example-8}
shows the overview tab for Nike Inc.\ over a thirty-day window.

\begin{figure}[H]
    \centering
    \includegraphics[width=\textwidth]{images/illustrative_example_8.png}
    \caption{Advertiser detail page for Nike Inc.}
    \label{fig:example-8}
\end{figure}

The \texttt{AdvertiserDetailPage} component issues two requests in parallel:
\texttt{GET /portal/advertisers/adv\_nike} returns the advertiser record
itself, and \texttt{GET /portal/advertisers/adv\_nike/analytics?days=30}
returns the rollup. The latter is computed in a single SQL query that
joins \texttt{events} to \texttt{ads} to \texttt{campaigns}, filters on
\texttt{advertiser\_id}, and aggregates impressions, clicks, engagements,
total budget, and total spent across every campaign owned by that
advertiser. A nested array of per-campaign summaries is returned in the
same payload to populate the campaigns tab.

The rollup is rendered as a row of eight KPI cards followed by a budget
utilisation bar that visualises \texttt{total\_spent} as a fraction of
\texttt{total\_budget}. The cards distinguish between counts (campaigns,
impressions, clicks, engagements), monetary values (budget, spend), and
derived ratios (CTR), giving the advertiser a single view of both
financial pacing and behavioural performance. The campaigns table below
uses status badges---green for active, amber for paused, blue for
completed, grey for draft---to surface lifecycle state without requiring
the operator to drill into each campaign.

% ============================================================

\subsection{Example 8: Campaign Detail --- Daily Pacing and Lifecycle Control}
\label{subsec:example-campaign-detail}

The third example narrows the lens further, to a single campaign.
Figure~\ref{fig:example-9} shows the Spring Running Campaign over a
thirty-day window, including the dual-axis daily breakdown chart and the
lifecycle action buttons.

\begin{figure}[H]
    \centering
    \includegraphics[width=\textwidth]{images/illustrative_example_9.png}
    \caption{Campaign detail page for the Spring Running Campaign.}
    \label{fig:example-9}
\end{figure}

The \texttt{CampaignDetailPage} component issues three requests:
\texttt{GET /portal/campaigns/\{id\}} for the campaign record,
\texttt{GET /portal/campaigns/\{id\}/analytics?days=30} for the metrics,
and \texttt{GET /portal/campaigns/\{id\}/ads?limit=100} for the
constituent ads. The analytics handler aggregates events scoped to ads
belonging to the campaign and additionally reads from the
\texttt{daily\_spend\_log} table, which is updated by the impression
logger on every served ad. The result is a \texttt{daily\_breakdown}
array containing one row per day with impressions, clicks, engagements,
and spend.

The breakdown is rendered as a Recharts dual-axis composite: bars for
daily impressions on the left axis, a line for daily spend on the right
axis, with a tooltip that surfaces all four metrics for the hovered day.
A budget progress bar and a today-spend indicator at the top of the page
provide instantaneous pacing feedback against the configured
\texttt{daily\_budget} cap. Lifecycle actions are exposed as buttons in
the header: \emph{Pause} issues \texttt{POST /portal/campaigns/\{id\}/pause},
which transitions the status to \texttt{paused} and immediately removes
the campaign's ads from the candidate pool in subsequent
\texttt{get\_candidate\_ads} calls; \emph{Resume} performs the inverse
transition.

% ============================================================

\subsection{Example 9: Engagement Analytics --- Beyond Click-Through Rate}
\label{subsec:example-engagement-analytics}

The fourth example demonstrates the engagement analytics surface, which
operationalises the multi-dimensional engagement signal discussed in
Section~\ref{subsec:adaptive-context}. Figure~\ref{fig:example-10} shows
the engagement page with one ad row expanded to reveal its sub-metrics.

\begin{figure}[H]
    \centering
    \includegraphics[width=\textwidth]{images/illustrative_example_10.png}
    \caption{Engagement analytics page showing the aggregate KPI strip,
    the engagement timeseries, and an expanded ad row with sub-metrics.}
    \label{fig:example-10}
\end{figure}

The \texttt{EngagementPage} component issues two requests:
\texttt{GET /analytics/engagement?days=7} returns per-ad aggregates
computed from the engagement events logged by the LLM-based engagement
classifier, and \texttt{GET /analytics/engagement/timeseries?days=7}
returns a daily series of engagement counts and average scores.
Aggregate KPIs across the displayed ads---total engagements, weighted
average engagement score, weighted average naturalness, positive and
negative sentiment counts---are computed client-side as engagement-weighted
means.

The page surfaces the central claim of the project's measurement
framework: a follow-up message that engages with an advertised product
constitutes a stronger signal of genuine interest than a click. The
expanded ad row in Figure~\ref{fig:example-10} illustrates this concretely.
Although Nike Pegasus 41 records a click-through rate of well under one
percent, its engagement rate exceeds fifteen percent and its average
naturalness score is 0.81, indicating that users frequently follow up on
the recommendation in conversational form even when they do not click
the tracking link. The colour coding on the sub-metrics rewards
naturalness and engagement scores above 0.7 and flags purchase proximity
below 0.4 as a candidate for context-level intervention.

% ============================================================

\subsection{Example 10: Adaptive Context Optimisation History}
\label{subsec:example-context-optimization}

The final example exposes the closed-loop mechanism that links the
engagement metrics of Example~\ref{subsec:example-engagement-analytics}
back to the prompt structure of Examples~\ref{subsec:example-product-research}
through~\ref{subsec:example-engagement}. Each ad in the inventory carries
a free-text \emph{presentation context guide} that is injected into the
response-synthesis prompt whenever the ad is matched. The guide is
versioned, and each successive version is generated by an LLM optimiser
that reviews the engagement metrics accumulated since the previous
version and proposes a revised guide together with a natural-language
rationale. The result is a per-ad audit trail of how the framing of each
product has evolved in response to behavioural feedback.

The history is surfaced inside the expanded ad row of the engagement
page and is fetched from \texttt{GET /admin/ads/\{ad\_id\}/context-history},
backed by the \texttt{ad\_context\_history} table. Each row of that table
carries the version number, the full \texttt{context\_text}, the
\texttt{optimization\_reasoning} authored by the optimiser LLM, the
\texttt{metrics\_snapshot} captured at the moment of optimisation, and a
timestamp. Manual optimisation can be triggered from the same panel via
\texttt{POST /admin/ads/\{ad\_id\}/optimize-context}; this call is
guarded by a minimum-engagement threshold of five and a minimum
inter-optimisation interval of twenty-four hours, both enforced server-side
in \texttt{optimize\_ad\_context}.

The three subfigures of Example 10 show the three
successive versions of the context guide for the Nike Pegasus 41 ad.
Version~1 (Figure~\ref{fig:example-11a}) is the initial guide generated
at ad-ingestion time from the structured ad fields; no metrics snapshot
is associated with it, and the reasoning field records only that the
guide was bootstrapped from the ad metadata.

\begin{figure}[H]
    \centering
    \includegraphics[width=0.8\textwidth]{images/illustrative_example_11a.png}
    \caption{Context guide version~1.}
    \label{fig:example-11a}
\end{figure}

Version~2 (Figure~\ref{fig:example-11b}) was produced after 180
impressions and 22 engagements had accumulated. The metrics snapshot
recorded at that moment shows an average engagement score of 0.55 and
an average purchase proximity of only 0.33. The optimiser's diagnostic
heuristics flagged the low purchase-proximity score as the dominant
deficiency and authored a revised guide that explicitly requires the
assistant to connect product features to the user's stated needs rather
than to recite generic benefits. The reasoning field records this
diagnosis verbatim.

\begin{figure}[H]
    \centering
    \includegraphics[width=0.7\textwidth]{images/illustrative_example_11b.png}
    \caption{Context guide version~2.}
    \label{fig:example-11b}
\end{figure}

Version~3 (Figure~\ref{fig:example-11c}) was produced after a further
65 impressions and 16 engagements. The snapshot at that moment shows
that purchase proximity had recovered to 0.41, but average naturalness
had fallen to 0.52---below the 0.6 threshold that triggers the
``soften the tone'' branch of the optimiser prompt. The revised guide
shifts the recommended register from assertive to peer-recommendation
style (``you might like'' rather than ``you should buy''), tightens the
relevance criterion to suppress the ad outside running-related queries,
and adds an explicit prohibition on price comparison. The reasoning
field credits the dismissal-rate signal as the secondary driver of the
relevance tightening.

\begin{figure}[H]
    \centering
    \includegraphics[width=0.8\textwidth]{images/illustrative_example_11c.png}
    \caption{Context guide version~3.}
    \label{fig:example-11c}
\end{figure}

The three versions together exhibit the central property of the
adaptive context mechanism: prompt-level adaptation against a behavioural
reward signal, with no gradient updates to the underlying model and a
fully interpretable change log. To the best of our knowledge, no prior
published work applies per-item LLM-driven prompt optimisation to a
conversational advertising setting.